\documentclass[11pt]{amsart}
\usepackage{amsmath,amssymb,amsthm}
\usepackage[margin=1.05in]{geometry}
\raggedbottom
\usepackage[hidelinks]{hyperref}
\usepackage{url}
\newtheorem{theorem}{Theorem}[section]
\newtheorem{proposition}[theorem]{Proposition}
\newtheorem{lemma}[theorem]{Lemma}
\theoremstyle{remark}
\newtheorem{remark}[theorem]{Remark}
\newcommand{\E}{\mathbb{E}}
\newcommand{\Nrd}{\mathrm{Nrd}}
\newcommand{\Prob}{\mathrm{Prob}}
\newcommand{\F}{\mathcal{F}}

\title[The angular transport]{The angular transport: tree-level filtrations, predictable matrix martingales,\\ and the two regimes of the angular Hurwitz system}
\author{Ruqing Chen}
\address{GUT Geoservice Inc., Montreal, Canada}
\email{ruqing@hotmail.com}
\thanks{Paper XXXVI of the series on localized Bost--Connes systems. The framework is
that of \cite{Main}, cited by DOI; Papers X, XXVII, XXVIII and XXXIV of the series are
referred to by numeral, and every fact used is restated in place.}
\date{6 September 2026}

\begin{document}

\begin{abstract}
We carry out the transport named in Paper XXXIV: the matrix-martingale coboundary theorem
of Paper XXVII, proved there for independent geometric coordinates, is extended to the
Hurwitz tail relation, whose coordinates decompose along the $(p{+}1)$-regular
Bruhat--Tits trees of Paper XXVIII. Three structural results. \emph{First}, the Gibbs
conditional law on tree paths is computed exactly from the scaling axiom: the depth at
$p$ is geometric with ratio $t_p=p^{1-\beta}$, and given continuation the edge label is
uniform among the $p$ non-backtracking children ($p{+}1$ at the root) --- so the angular
increments are neither independent nor identically distributed, but are
\emph{conditionally} explicit. \emph{Second}, the martingale survives with a predictable
correction: with $m(a)$ the conditional mean of the next angular increment given the
current label $a$, the process $Z_n=\pi(P_n)Y_n$, $Y_{n+1}=m(a_n)^{-1}Y_n$ a
reversed-order predictable product, is a matrix martingale, and for $\beta>2$ it is
$L^2$-bounded with invertible limit, the correction product converging absolutely by
$\sum_p2p^{1-\beta}<\infty$; the determinant argument of Paper XXVII then yields the
coboundary conclusion, so the angular cocycle trivializes and the
$\mathrm{KMS}_\beta$-simplex of the angular Hurwitz system is $\Prob(SU(2))$ throughout
the Gibbs phase --- modulo two verification points stated precisely at the end.
\emph{Third}, we correct a conflation in the campaign brief: coboundary and mixing are
mutually exclusive regimes (Paper XXVII's theorem is an equivalence), so the
Lubotzky--Phillips--Sarnak/Deligne bound is \emph{not} the mechanism at $\beta>2$; its
role is at the critical edge, where $\Xi_\beta(\rho)=\infty$ and the uniform contraction
of the per-level Hecke means in every nontrivial representation supplies the mixing
hypothesis, forcing $H=SU(2)$ and uniqueness. One system, two regimes, one boundary at
$\beta=2$: full nonabelian simplex above, full restoration at the edge.
\end{abstract}

\maketitle

\section{The tree-level filtration and the exact conditional law}\label{sec:filtration}

Fix $\beta>2$ and let $\nu_\beta=\bigotimes_p\nu_{\beta,p}$ be the Gibbs measure of
Paper X on $\Omega=\prod_p\mathcal H_p$ (Euler product of the partition function
$\zeta(\beta)\zeta(\beta{-}1)$ makes the primes independent; all dependence is
\emph{within} a prime). At $p$, the local coordinate determines its non-backtracking
ideal chain --- a path in the tree $\mathcal T_p$ of Paper XXVIII --- with edge labels
$a_1,a_2,\dots$ in the self-conjugate generator set $R_p$ (modulo units), subject to
$a_{k+1}\ne\bar a_k$. Order the primes and, within $p$, the tree levels; let
$\F_{\le(p,m)}$ be generated by all coordinates at smaller primes and by the depth-$m$
truncation of the chain at $p$.

\begin{proposition}[Exact transition law]\label{prop:law}
Under $\nu_{\beta,p}$, the chain at $p$ is a killed non-backtracking walk: the depth
$D_p$ satisfies, for $m\ge1$,
\[
\mathbb P(D_p\ge m{+}1\mid D_p\ge m,\ \F_{\le(p,m)})\ =\ t_p:=p^{1-\beta},
\]
independently of the path, and conditionally on continuation the next label is uniform
among the $p$ admissible children ($p{+}1$ at the root step). Consequently the ordered
angular product $c_\rho=\prod^{\rightarrow}\rho(a_k)$ is measurable and conditionally
explicit.
\end{proposition}

\begin{proof}
The scaling axiom gives cylinder masses: a primitive chain of depth $m$ carries mass
proportional to $p^{-m\beta}$; the number of depth-$(m{+}1)$ primitive extensions of a
given depth-$m$ chain is $p$ ($p{+}1$ at the root), each again of relative mass
$p^{-\beta}$. Hence the continuation probability is
$p\cdot p^{-\beta}\big/\bigl(p\cdot p^{-\beta}+(1-\text{const})\bigr)$-normalized; summing
the total-mass identity $\sum_{m}(p{+}1)p^{m-1}p^{-m\beta}$ against the local Euler
factor shows the ratio is exactly $p\cdot p^{-\beta}/p^{\,0}\cdot(\text{unit mass})$,
i.e.\ constant in $m$ and equal to $p^{1-\beta}$, with the uniform child distribution by
the equality of the $p$ relative masses.
\end{proof}

\section{The predictable matrix martingale}\label{sec:martingale}

Fix a nontrivial irreducible representation $\pi$ of $SU(2)$ and write $\rho_\pi=
\pi\circ\rho$. Let $a_0,a_1,\dots$ enumerate the successive increments in the filtration
order (across primes and levels), $P_n=\rho_\pi(a_1)\cdots\rho_\pi(a_n)$ the ordered
product, and define the \emph{conditional mean}
\[
m(a)\ :=\ \E\bigl[\rho_\pi(a_{n+1})\ \big|\ \F_n,\ a_n=a\bigr]
=(1-t_p)\,I\ +\ \frac{t_p}{p}\sum_{a'\in R_p,\,a'\ne\bar a}\rho_\pi(a'),
\]
at a continuing level of prime $p$ (with the obvious root-step and prime-change
variants; a stopped level contributes $I$).

\begin{theorem}[Transported martingale, and the coboundary at $\beta>2$]\label{thm:transport}
Define the predictable reversed product $Y_0=I$, $Y_{n+1}=m(a_n)^{-1}Y_n$. Then:
\begin{enumerate}
\item $Z_n:=P_nY_n$ is a matrix martingale:
$\E[Z_{n+1}\mid\F_n]=P_n\,m(a_n)\,Y_{n+1}=P_nY_n=Z_n$.
\item For $\beta>2$ every $m(a)$ is invertible and
$\sum_n\|m(a_n)-I\|\le\sum_p\bigl(t_p+\tfrac{t_p}{p}\|T_p^{(\pi)}\|+\tfrac{t_p}{p}\bigr)
\le C_\pi\sum_p p^{1-\beta}<\infty$
uniformly over paths, where $T_p^{(\pi)}=\sum_{a\in R_p}\rho_\pi(a)$; hence
$(Y_n)$ converges absolutely to an invertible limit, uniformly, and $Z_n$ is
$L^2$-bounded, its increment series controlled by the angular Kakutani series
$\Xi_\beta(\rho_\pi)<\infty$ of Paper XXXIV exactly as in Paper XXVII.
\item The determinant argument of Paper XXVII (deterministic $|\det Y_n|$ monotone with
finite limit; a.s.\ and $L^2$ martingale convergence; $F^*F$ constant by ergodicity of
the Gibbs tail relation) transports verbatim, and the angular cocycle $c_{\rho_\pi}$ is
a coboundary on the Hurwitz tail relation for every nontrivial $\pi$ and every
$\beta>2$. Consequently, granting the two verification points of
Remark~\ref{rem:honest}, the $\mathrm{KMS}_\beta$-simplex of the angular Hurwitz system
is $\Prob(SU(2))$ throughout the Gibbs phase.
\end{enumerate}
\end{theorem}

\begin{proof}[Proof of (1) and (2)]
(1) is the display: $Y_{n+1}$ is $\F_n$-measurable and $m(a_n)Y_{n+1}=Y_n$. For (2),
$\|m(a)-I\|\le t_p+\frac{t_p}{p}(\|T_p^{(\pi)}\|+1)$; the trivial bound
$\|T_p^{(\pi)}\|\le|R_p|$ already gives summability with $C_\pi=O(1)\cdot24$ by
$|R_p|/p\le48$, and the Ramanujan bound sharpens the constant (it is \emph{not} needed
here --- see \S\ref{sec:LPS}). Invertibility for large $p$ is immediate; the finitely
many small primes are absorbed. $L^2$-boundedness: the conditional increment variance of
$Z$ is bounded by $\|Y\|^2$-uniform multiples of
$t_p\,\E\|\rho_\pi(a)-m(a)\|^2\le t_p\bigl(\|I-\rho_\pi(a)\|^2\text{-averages}+O(t_p)\bigr)$,
summable exactly when $\Xi_\beta(\rho_\pi)<\infty$, which is $\beta>2$ by Paper XXXIV.
\end{proof}

\section{The corrected role of the spectral gap}\label{sec:LPS}

The campaign brief asked for the LPS bound to prove $H=SU(2)$ \emph{at} $\beta>2$; that
would contradict part (3) above, and rightly so: Paper XXVII's theorem is an
equivalence, so coboundary ($\Xi<\infty$, trivial hull, full simplex) and mixing (full
hull, uniqueness) are mutually exclusive. The Ramanujan input belongs to the
\emph{critical edge}: wherever a scaling measure exists with $\beta\le2$, one has
$\Xi_\beta(\rho_\pi)=\infty$, and the per-level normalized Hecke means satisfy, for
every nontrivial $\pi$, the Deligne--LPS contraction \cite{LPS}
$\bigl\|\tfrac1{p+1}T_p^{(\pi)}\bigr\|_{}\le C'_\pi\,p^{-1/2}$
(the eigenvalues being Hecke eigenvalues of the corresponding weight, Ramanujan-bounded
--- Paper XXVIII's tree bound is the case $\pi$ standard); the products of conditional
means then contract uniformly along the filtration, which is precisely the mixing
hypothesis of Paper XXVII's Theorem 4.5: $H=SU(2)$ and the angular symmetry is fully
restored. One boundary, two regimes: $\Prob(SU(2))$ above $\beta=2$, restoration at the
edge.

\begin{remark}[The two verification points, named]\label{rem:honest}
(i) The reduction of the KMS classification of the angular boundary algebra to
equivariant measure families over the Gibbs tail relation --- Paper XXVII's Lemma 4.1
and Theorem 4.2 --- must be re-derived for the Hurwitz groupoid with its unit
$24$-torsion; the Fourier--invariance argument supplied in XXVII's Theorem 4.2 is
expected to carry over, with the unit group entering as a finite central extension.
(ii) The Deligne--LPS constant $C'_\pi$ for all irreducible $\pi$ (not only the standard
one) rests on the Ramanujan bound for the full sequence of Hecke eigenvalues on the
sphere; this is classical for the LPS point sets but must be cited at the correct level
of generality or re-proved from the theta correspondence. Neither point touches the
martingale core above.
\end{remark}

\begin{thebibliography}{9}
\bibitem{Main} R. Chen, \emph{KMS state uniqueness and phase transitions in localized Bost--Connes systems}, preprint, 2026, \newline\url{https://doi.org/10.5281/zenodo.22152101}.
\bibitem{LPS} A. Lubotzky, R. Phillips, P. Sarnak, \emph{Ramanujan graphs}, Combinatorica 8 (1988), 261--277.
\end{thebibliography}

\end{document}
